% !TEX root = ../main.tex

\section{Metodolog\'ia}
\label{sec:methodology}

\subsection{Dise\~no del lenguaje}
\label{sec:methodology:design}

\sculpt es, en rigor, dos lenguajes acoplados. El primero es un
\emph{lenguaje visual y f\'isico}: una especificaci\'on de los
componentes que pueden existir y de las reglas de construcci\'on entre ellos.
El segundo es un \emph{lenguaje de programaci\'on} que define el
comportamiento que cada componente representa cuando se ensambla
junto a otros. Cada programa \sculpt es, simult\'aneamente, un objeto
construible en el espacio tridimensional y una expresi\'on computable.

El dise\~no del lenguaje sigui\'o un enfoque centrado en el usuario,
con el objetivo expl\'icito de que el conjunto de componentes fuera
reconocible para audiencias diversas. Cada componente se distingue
mediante una forma simple pero \'unica, de modo que pueda
identificarse sin recurrir a conocimiento previo ni a la lectura de
texto. La sintaxis del lenguaje es una propiedad
geom\'etrica: dos piezas encajan o no encajan, y la legalidad del
programa se inspecciona f\'isica y mecanicamente.

El conjunto de operaciones es deliberadamente m\'inimo. \sculpt tiene
solo 13 instrucciones que operan sobre pilas FILO de n\'umeros
racionales. Las pilas son la \'unica estructura de datos del lenguaje,
y se crean impl\'icitamente la primera vez que una instrucci\'on las
referencia. Cualquier forma no reservada para una instrucci\'on puede
designar una pila, lo que extiende la libertad est\'etica del lenguaje
tambi\'en al espacio de datos. Sobre este conjunto m\'inimo se obtiene
Turing-completitud, como se demuestra en la seccion \ref{sec:results:turing}.

La est\'etica de los componentes es, por dise\~no, abierta. La
especificaci\'on solo restringe las conexiones entre piezas y el
conjunto de \emph{features} que cada bloque debe ofrecer; la forma
final, el tama\~no, el material, el color y la textura quedan al
arbitrio de quien construye el programa. Un cubo PLA est\'andar, una
figura tallada en madera o un trozo de papel garabateado y mal doblado 
son piezas v\'alidas mientras los puntos de
acople respeten la especificaci\'on. Esta apertura es lo que permite,
al mismo tiempo, ajustar las piezas a necesidades de accesibilidad
particulares (tama\~no, peso, contraste, textura) y abrir el lenguaje
como medio de expresi\'on art\'istica.

\subsection{Implementaci\'on f\'isica y \emph{runtime}}
\label{sec:methodology:implementation}

La implementaci\'on de referencia utilizada en las pruebas consta de
piezas peque\~nas impresas en \'acido polil\'actico (PLA) de
colores variados y piezas de melamina blanca, conectadas mediante imanes y juntas roscadas tipo
codo. Esta combinaci\'on permite el ensamblaje y desensamblaje
repetido sin herramientas, resistencia esctructural y facilita modificar el programa durante
la sesi\'on.

La fisicalidad del lenguaje conlleva un compromiso. Al ser un medio
estrictamente f\'isico, la ejecuci\'on de un programa \sculpt depende
de un \emph{int\'erprete humano}, que recorre las piezas y aplica las
instrucciones. La interpretaci\'on humana es parte del ejercicio de
\sculpt como acto performativo, pero no es confiable. Los errores
son inevitables y el lenguaje no es tolerante a fallos en caso de errores de interpretación.
Para compensar esta fragilidad, el ecosistema
incluye \sculpter (\emph{Simple Cubic Language for Programming
Task's Environment and Runtime}), un transpilador escrito en Scala~3
que toma la representaci\'on textual de un programa \sculpt y lo
ejecuta paso a paso. Gracias a Scala.js, \sculpter corre en cualquier
navegador moderno, sin instalaci\'on adicional, lo que lo hace
distribuible y accesible para quien quiera probar un programa antes
de ensamblarlo f\'isicamente.

\subsection{Protocolo de validaci\'on}
\label{sec:methodology:validation}

La validaci\'on se realiz\'o en seis sesiones con
participantes repartidos en perfiles diversos: novatos sin
experiencia previa en programaci\'on, programadores expertos con
a\~nos de experiencia, estudiantes de secundaria, ni\~nos y
profesionales del campo art\'istico. La diversidad de perfiles fue
intencional. \sculpt aspira a ser inclusivo, y la validaci\'on debe
ponerlo a prueba con pre\'ambulos cognitivos y motivacionales muy
distintos.

Cada sesi\'on sigui\'o cuatro fases. (1) Una \emph{introducci\'on}
de 30 a 45~minutos, adaptada al perfil del grupo: m\'as detallada
sobre los componentes f\'isicos para los novatos y m\'as t\'ecnica
para los expertos. (2) Una fase de \emph{exploraci\'on libre} de
una hora, sin tareas asignadas, en la que los participantes
manipulan las piezas a su ritmo para familiarizarse con sus formas
y conexiones. (3) Una fase de \emph{tareas computacionales} de
dos horas, con conjuntos de problemas calibrados al nivel de cada
grupo (\eg sumas iteradas, condicionales para romper bucles); los
facilitadores resolv\'ian dudas pero el participante completaba la
tarea por su cuenta. (4) Una \emph{encuesta y entrevista} final opcional
sobre usabilidad, accesibilidad y experiencia general; en el caso
de menores de edad, la encuesta tambi\'en la respondieron los
docentes acompa\~nantes, con preguntas adicionales sobre la
viabilidad de \sculpt como herramienta pedag\'ogica.

El an\'alisis combin\'o m\'etodos cualitativos y cuantitativos.
En lo cuantitativo se midi\'o el tiempo a tarea y el n\'umero de
errores cometidos durante la fase de programaci\'on. En lo
cualitativo se analizaron las entrevistas y las preguntas abiertas
de la encuesta, buscando patrones recurrentes en la percepci\'on
de las piezas, en los obst\'aculos sint\'acticos y en la
experiencia est\'etica de manipular el lenguaje. Adicionalmente,
con el subgrupo de artistas se realizaron sesiones extra
dedicadas exclusivamente a la creaci\'on de piezas no funcionales,
con el fin de evaluar a \sculpt como medio expresivo independiente
de su capacidad computacional.

\endinput
